% Fate's Edge SRD −−− Complications & Timers (Section 18)

\section{Complications & Timers}

Complications are the narrative, mechanical, and environmental consequences that arise from failure, pressure, or the GM spending Story Beats. Timers are the primary tool for tracking accumulating threats or progress.

\subsection{Story Beats (SB)}

Story Beats are the GM's currency for introducing complications. Each die result of 1 generates 1 SB. Re−rolling a 1 does not erase its SB; if the re−rolled die also shows 1, it generates additional SB.

⫕section{SB Spend Menu}

\begin{tabular}{@{}ll@{}}
⊤rule
\textbf{SB Cost} & \textbf{Effect} \\
∣rule
1 & Minor pressure: noise, trace, tick a 4–6 segment timer +1, worsen Position by one step. \\
2 & Moderate setback: alarm raised, lose cover/advantage, a lesser foe appears, split the party briefly. \\
3 & Serious trouble: reinforcements arrive, key gear breaks (temporarily), a major twist shifts the scene. \\
4+ & Scene−shaping turn: a trap springs, authority arrives, the environment changes dramatically, a Patron manifests an omen. \\
⊥tomrule
\end{tabular}

⫕section{Limits on SB per Scene}
\begin{itemize}
  ℹtem Base SB budget: 4 + highest character Tier (e.g., Tier I = 5, Tier III = 7).
  ℹtem Maximum 12 SB for standard scenes, 16 SB for climactic scenes.
  ℹtem Unspent SB may carry over between scenes (GM discretion, max 4).
\end{itemize}

\subsection{Boons (Player Resource)}

Boons are the players' resource for mitigating complications or turning the tide. See Section~\ref{sec:core−mechanic} for detailed earning and spending rules.

\subsection{Timers}

Timers track accumulating threats or progress. They are a visible way to manage pressure.

⫕section{When to Use a Timer }
\begin{itemize}
  ℹtem \textbf{Threat Timer } (4–8 segments): Fills with enemy action, environmental decay, or social collapse. When full, a defined disaster occurs.
  ℹtem \textbf{Progress Timer } (4–8 segments): Fills with player successes to achieve a goal (e.g., escape, research, negotiation).
  ℹtem \textbf{Contest Timers:} Two timers race – first to fill wins.
\end{itemize}

⫕section{Ticking Timers}
Timers advance when:
\begin{itemize}
  ℹtem The GM spends SB (e.g., 1 SB = tick a relevant timer +1).
  ℹtem A failed roll directly relates to the threat (GM discretion).
  ℹtem A scene transition or scripted event occurs.
\end{itemize}

⫕section{Example Timers}
\begin{itemize}
  ℹtem \textbf{Guards Incoming [6]} – ticks on noise or missed stealth. On fill, a patrol arrives.
  ℹtem \textbf{Collapsing Bridge [4]} – ticks on failed Athletics or heavy impacts. On fill, bridge falls.
  ℹtem \textbf{Patron's Patience [8]} – ticks each time Obligation is ignored. On fill, Patron demands service.
\end{itemize}

\subsection{Hazards}

Hazards are ongoing complications that require action to avoid or mitigate.

⫕section{Environmental Hazards}
\begin{itemize}
  ℹtem \textbf{Active:} Fire spreading, flooding, toxic gas – represented by a Threat Timer that ticks each round or on SB spends.
  ℹtem \textbf{Passive:} Difficult terrain, darkness, extreme heat/cold – impose Position penalties or DV modifiers.
\end{itemize}

⫕section{Social Hazards}
\begin{itemize}
  ℹtem \textbf{Rumors:} A timer that advances when the party fails social rolls; when full, a faction turns hostile.
  ℹtem \textbf{Scrutiny:} Worsens Position for stealth or deception.
  ℹtem \textbf{Reputation Damage:} Future social rolls start at worse Position.
\end{itemize}

⫕section{Magical Hazards}
\begin{itemize}
  ℹtem \textbf{Backlash Zone:} An area where free casting automatically generates +1 SB per spell.
  ℹtem \textbf{Ward / Seal:} Requires a roll to pass; failure triggers a consequence (harm, alarm, etc.).
  ℹtem \textbf{Leaking Veil:} Spirits or fey occasionally manifest spontaneously (GM spends SB to introduce).
\end{itemize}

\subsection{Fear Effects}

When a character escalates on the Fear Track (Shaken $→$ Frightened $→$ Panicked), roll or choose an effect from the table below. Apply to NPCs by default; PCs may adopt as narrative guidance.

\begin{tabular}{@{}rll@{}}
⊤rule
\textbf{d10} & \textbf{Effect} & \textbf{Magic Tags} \\
∣rule
1 & Freeze: Cannot act this round, staring or trembling. & Silence, Stasis \\
2 & Flee: Must move at full speed away from the source of Fear. & Movement, Wind \\
3 & Drop: Drops what they are holding. & Disarm, Break \\
4 & Beg: Pleads or bargains incoherently. & Compulsion, Voice \\
5 & Hide: Seeks cover, concealment, or allies. & Shadow, Illusion \\
6 & Attack in Panic: Lashes out wildly at the nearest target. & Rage, Fire \\
7 & Blunder: Stumbles into danger (trap, hazard, off balance). & Chaos, Trickery \\
8 & Obey: Instinctively follows a simple command from the fear−causer. & Command, Charm \\
9 & Break Down: Sobs, prays, or becomes useless until aided. & Curse, Despair \\
10 & Catatonia: Becomes unresponsive, requiring intervention. & Sleep, Dream \\
⊥tomrule
\end{tabular}

⫕section{Fear Track Escalation}
\begin{itemize}
  ℹtem \textbf{Shaken:} Apply minor effects (hesitation, –1 die, startled).
  ℹtem \textbf{Frightened:} Roll normally on the table.
  ℹtem \textbf{Panicked:} Apply severe or exaggerated results (e.g., 2 = reckless flight into danger, 6 = attacks allies).
\end{itemize}

\subsection{The Deck of Consequences (Optional)}

The Deck of Consequences is a narrative tool that turns Story Beats into thematic, unpredictable complications. It replaces or supplements GM fiat with randomized twists.

⫕section{Deck Structure}
Use a standard 52−card deck (jokers optional). Suits map to complication themes:

\begin{tabular}{@{}ll@{}}
⊤rule
\textbf{Suit} & \textbf{Theme} \\
∣rule
Hearts   & Social / Emotional \\
Spades   & Physical / Harm \\
Clubs    & Material / Cost \\
Diamonds & Magical / Spiritual \\
⊥tomrule
\end{tabular}

⫕section{Rank Severity}
\begin{itemize}
  ℹtem 2–5: Minor inconvenience (noise, Fatigue 1, small delay).
  ℹtem 6–10: Moderate setback (alarm raised, gear broken, Position worsened).
  ℹtem J, Q, K: Major consequence (reinforcements, key gear destroyed, front timer advanced).
  ℹtem A: Severe twist (trap springs, authority arrives, scene shifts).
\end{itemize}

⫕section{When to Draw}
The GM may draw from the deck when:
\begin{itemize}
  ℹtem A roll shows two or more 1s.
  ℹtem The GM spends 2+ Story Beats at once.
  ℹtem A major timer fills (scene or campaign timer).
  ℹtem During travel to generate an unexpected encounter.
\end{itemize}
Aim for 1–2 draws per session; more if the tone skews chaotic.

⫕section{Procedure}
\begin{enumerate}
  ℹtem Note the number of SB spent (or generated).
  ℹtem Draw a number of cards equal to the SB spent (minimum 1, maximum 3).
  ℹtem Synthesize a single narrative complication from the combination of suits and the highest rank.
  ℹtem Do not reshuffle until the deck runs out or the session ends – this creates a natural ``fate budget.''
\end{enumerate}

⫕section{Shortcut Table for Synthesis}
\begin{tabular}{@{}ll@{}}
⊤rule
\textbf{Suit} & \textbf{Low (2–5)} & \textbf{High (J–A)} \\
∣rule
Hearts   & Awkward social moment & Betrayal, scandal, broken trust \\
Spades   & Bruise, small stumble & Serious injury, structural collapse \\
Clubs    & Minor loss (rations, coins) & Major asset destroyed, debt called \\
Diamonds & Omen, eerie noise & Patron demand, reality bends \\
⊥tomrule
\end{tabular}

\subsection{Integrating Complications with SB and Boons}

\begin{itemize}
  ℹtem SB and Boons create a push−pull: SB gives the GM tools to escalate; Boons give players tools to recover or resist.
  ℹtem A well−spent SB should change the situation, not just punish. Offer the player a Devil's Bargain: ``I will spend 2 SB to have the guard turn around – but you will owe a favor.''
  ℹtem Timers make threat visible; when a timer ticks, describe what changes in the fiction.
\end{itemize}

\subsection{Example: Complication in Action}

Lyra misses a stealth roll (0 successes) while sneaking past a guard post. She gains 2 Boons. The GM spends 1 SB: ``Your foot scrapes a loose stone – the guard looks up.'' (Tick \textit{Guards Alert [4]} +1.)

Later, the GM spends another 2 SB: ``The guard calls for backup – two more arrive. New timer: \textit{Reinforcements [6]}.'' Lyra uses 1 Boon to re−roll a failed athletics test to climb a wall, escaping just as the timer reaches 4.

Another scene: Lyra rolls three 1s on a critical casting. The GM draws three cards from the Deck of Consequences: Knave of Hearts, 7 of Clubs, Ace of Spades. Synthesis: ``You slip past the first guard, but your cloak catches on a rusty nail (resource drain). Worse, the second guard recognizes your face from a wanted poster – she is your cousin (social betrayal). She hesitates, then shouts `Traitor!' As she draws her blade, a hidden portcullis drops behind you (scene shift). You are trapped in the courtyard.''

